\section{Literature Review}
\label{sec:literature}

\subsection{Existing Research}

Other researchers have already studied parts of this problem, and their work forms the starting point for this thesis.

First, studies have shown that AI can read and extract information from messy, unstructured documents like scanned forms and receipts. Mathew et al.\ \cite{mathew_docvqa_2021} created a large public dataset called DocVQA, which is one of the most widely used tools for testing how well an AI can answer questions from a document image. Jaume et al.\ \cite{jaume2019funsd} created a well-known test collection of scanned administrative forms, and Park et al.\ \cite{park2019cord} did similar work using a collection of receipts. These studies all show that extracting data correctly from real-world documents is genuinely difficult, and that a proper test dataset is needed to measure it fairly.

Second, researchers have shown that large AI models can be made much smaller without losing much of their ability. Dettmers et al.\ \cite{dettmers2022llmint8} showed that a large model can be compressed to use much less memory and still work well on a regular laptop. Frantar et al.\ \cite{frantar2022gptq} went even further with a method that compresses models even more efficiently. A free and open-source tool called Ollama \cite{ollama_2024} now makes it straightforward to download and run these smaller models on any consumer laptop --- the tool handles all setup and hardware compatibility automatically. Three specific models are used in this study: \texttt{deepseek-ocr} by DeepSeek AI \cite{deepseek_ocr_2024}, a document-focused vision model; \texttt{gemma4:e4b} by Google DeepMind \cite{gemma4_2026}, a multimodal model that reads document images; and \texttt{qwen3.5:4b} by Alibaba \cite{qwen3_2025}, a model strong at structured extraction. A comparison of recent models in this class shows that they each perform differently on different types of documents \cite{gemma4_phi4_qwen3_2026}. Taken together, this body of work proves that running a capable, document-aware AI on a normal laptop is now possible --- and that the choice of which local model to use can matter a lot.

Third, researchers have studied how to measure AI accuracy fairly. Peer et al.\ \cite{anlsstar_2024} introduced a metric called ANLS* that is much better suited to generative AI models than older yes/no scoring methods. It gives partial credit for almost-correct answers and can check structured outputs like JSON field by field. Zhang et al.\ \cite{zhang_bertscore_2020} introduced BERTScore, which checks whether the meaning of an AI's answer matches the correct answer, even if the exact words are different. These newer metrics are more honest measures of how well an AI is really performing.

Fourth, Strubell et al.\ \cite{strubell2019energy} studied how much energy and money it costs to run AI systems and provided a clear method for calculating those costs. Wei et al.\ \cite{wei2022emergent} also studied the general abilities and limits of large AI models, which helps explain why some models perform better than others on specific tasks.

\subsection{Summary of Key Related Work (2024--2026)}

Table~\ref{tab:related-work} gives a quick overview of the most directly related studies from the last two years, showing what each one found, how it was done, and what question it left unanswered.

\begin{table}[h!]
\centering
\small
\caption{Summary of Related Work (2024--2026)}
\label{tab:related-work}
\begin{tabular}{|p{0.7cm}|p{2.5cm}|p{3.7cm}|p{3.2cm}|p{3.7cm}|}
\hline
\textbf{Year} & \textbf{Study} & \textbf{Key Finding} & \textbf{Methodology} & \textbf{Gap Left} \\
\hline
2024 &
Peer et al.\ \cite{anlsstar_2024} \newline (ANLS*) &
Introduced a metric that fairly scores generative LLM document extraction, giving partial credit for near-correct answers &
Extended ANLS with hierarchical JSON tree matching using the Hungarian algorithm &
Did not compare local vs.\ cloud models or include any cost or latency analysis \\
\hline
2024 &
TokenPowerBench \cite{tokenpowerbench_2024} &
Provided a reproducible method to measure energy used per LLM inference call on real hardware &
Software-level power profiling using \texttt{powermetrics} and NVML for GPU wattage &
No document extraction task; no cost-vs-accuracy comparison provided \\
\hline
2024 &
Shirazi et al.\ \cite{arial_2024} \newline (ARIAL) &
Agentic pipeline for document VQA with exact pixel-level answer localization, beating prior end-to-end models &
Modular OCR + RAG retrieval + fine-tuned SLM for final answer generation &
Complex system not suited for simple enterprise deployment; no operational cost model \\
\hline
2025 &
Microsoft Research \cite{phi4mini_2025} \newline (Phi-4 Mini) &
A 3.8B model achieves strong document understanding using a Mixture of LoRAs and a multi-crop vision strategy &
Dynamic multi-crop vision (up to 36 crops) + DocVQA and benchmark evaluation &
No comparison against cloud APIs; no cost or latency analysis included \\
\hline
2025 &
Qwen Team \cite{qwen25vl_2025} \newline (Qwen2.5-VL) &
Achieves state-of-the-art structured Key Information Extraction directly from document images &
Dynamic resolution training with absolute spatial coordinates on diverse document data &
Single model only; no local vs.\ cloud comparison; no cost model provided \\
\hline
2026 &
Google DeepMind \cite{gemma4_2026} \newline (Gemma 4) &
Runs multimodal AI efficiently on edge devices with only 4.5B active parameters per forward pass &
Per-Layer Embeddings + alternating attention design + open multimodal benchmark evaluation &
No cost modeling; no head-to-head test against cloud services or other local models \\
\hline
2026 &
Lenovo Press \cite{lenovo_tco_2026} \newline (TCO Report) &
On-premises AI reaches financial breakeven vs.\ cloud in under 4 months at 20\%+ hardware utilization &
CapEx/OpEx modeling with straight-line depreciation and energy cost analysis (PUE formula) &
No accuracy metric; no specific model tested; no document extraction task evaluated \\
\hline
\end{tabular}
\end{table}

\subsection{Research Gap}

Even though all of these areas have been studied separately, nobody has put them all together in one place. In particular, no study has compared several local AI models --- each built in a different way and run through a common tool like Ollama --- against paid online services like GPT-4o \cite{openai_gpt4o_2024} and Gemini 3.1 Pro \cite{google_gemini_2025} on the same task, measuring accuracy, speed, and cost at the same time. Testing only one local model is a problem because if it performs badly, it is impossible to tell whether that is a weakness of that one model or a weakness of local AI tools in general \cite{gemma4_phi4_qwen3_2026}. At the same time, no study has checked whether accuracy results change depending on how the instruction to the AI is written, even though this is a well-known issue with AI models \cite{anlsstar_2024}. This thesis fills both of these gaps.

\subsection{Research Questions}

Based on everything studied above, this thesis will answer these three specific questions:

\begin{itemize}
    \item \textbf{Question 1:} Is there a real, meaningful difference in accuracy between the three local AI models and the paid online service when extracting information from the same documents, and do those differences stay the same across different ways of writing the instruction?
    \item \textbf{Question 2:} How much faster or slower is each local AI model compared to the paid online service?
    \item \textbf{Question 3:} At what number of documents per day does it become cheaper to use each local AI model instead of the paid online service, based on a clear and repeatable cost calculation?
\end{itemize}
